\documentclass[11pt,a4paper]{article}

% ------------------------------------------------------------------
% Packages
% ------------------------------------------------------------------
\usepackage[T1]{fontenc}
\usepackage[utf8]{inputenc}
\usepackage{lmodern}
\usepackage{microtype}
\usepackage[a4paper,margin=1in]{geometry}
\usepackage{setspace}
\setstretch{1.08}

\usepackage{amsmath,amssymb,amsfonts}
\usepackage{graphicx}
\usepackage{booktabs}
\usepackage{array}
\usepackage{tabularx}
\usepackage{longtable}
\usepackage{enumitem}
\usepackage[font=small,labelfont=bf,labelsep=period]{caption}
\usepackage{algorithm}
\usepackage{algpseudocode}
\usepackage{xcolor}
\usepackage{fancyhdr}
\usepackage{cite}
\usepackage{hyperref} % load last

\hypersetup{
    colorlinks=true,
    linkcolor=blue!60!black,
    citecolor=blue!60!black,
    urlcolor=blue!60!black,
    pdftitle={Multi-Slice Privacy-Aware Traffic Forecasting at RAN Level},
    pdfauthor={Undergraduate Research Team (Computer Networks)}
}

\setlist{itemsep=2pt,topsep=4pt}
\renewcommand{\arraystretch}{1.2}
\newcolumntype{Y}{>{\raggedright\arraybackslash}X}

% Header / footer
\pagestyle{fancy}
\fancyhf{}
\lhead{\small\itshape Multi-Slice Privacy-Aware Traffic Forecasting at RAN Level}
\rhead{\small\thepage}
\renewcommand{\headrulewidth}{0.4pt}

% ------------------------------------------------------------------
% Title block
% ------------------------------------------------------------------
\title{\textbf{Multi-Slice Privacy-Aware Traffic Forecasting at RAN Level:\\[2pt]
A Rigorous Replication, Critical Audit, and Resilient Extensions}}

\author{
    \textbf{Undergraduate Research Team --- Computer Networks}\\[6pt]
    \small
    \begin{tabular}{@{}llp{8.2cm}@{}}
        Ankush R            & CS24B1062 & Paper Analyst: problem formulation, system model and architectural scrutiny\\
        Vihaan Gupta        & CS24B1063 & Reproduction Lead: data pipeline, baselines and federated core implementation\\
        Piyush Mishra       & CS24B1048 & Scalability Lead: clustered FL replication, algorithmic debugging and communications\\
        Narayana Murthy S   & CS24B1047 & Extension Lead: critical audit, privacy bounds and empirical improvements\\
    \end{tabular}\\[6pt]
    \small Repository: \url{https://github.com/pengeon1/CN_Research}
}
\date{}

\begin{document}

\maketitle
\thispagestyle{plain}

% ------------------------------------------------------------------
% Abstract
% ------------------------------------------------------------------
\begin{abstract}
\noindent
Predictive radio resource allocation across heterogeneous 5G and next-generation slices requires fine-grained traffic forecasting while respecting strict multi-tenant privacy boundaries between Infrastructure Providers (InPs) and Mobile Virtual Network Operators (MVNOs). This report presents an independent replication, critical architectural audit, and empirical extension of the federated learning framework proposed by Phyu, Stanica, and Naboulsi (\textit{IEEE TNSM}, Dec.\ 2023). We re-implement Centralized Long Short-Term Memory (LSTM), Decentralized local LSTM (DLSTM), Federated Averaging (FALSTM), FedProx-based FL (FPLSTM), and Clustered FL schemes (IC-FPLSTM and RC-FPLSTM) on a calibrated metropolitan cellular environment spanning 57 base stations and 4 application slices over 10 days at 10-minute granularity. Our multi-seed evaluation shows that federated learning reduces the local DLSTM forecasting error by 14.3\% (RMSE 0.1748 vs.\ 0.2039). The audit also identifies five principal flaws in the original formulation: (1) a divergence between the proximal objective and the post-hoc descent of Algorithm~2; (2) an unweighted arithmetic aggregation in Algorithm~1 that contradicts the paper's own sample-weighted objective, Eq.~(1); (3) a dynamic set-mutation bug in Algorithm~4 that terminates clustering prematurely and skews cluster balance by up to 33\%; (4) an overstated clustering benefit, with only a 1.24\% improvement over random clustering; and (5) a privacy paradox in IC-FPLSTM, which requires sharing raw 24-hour traffic vectors with the coordinator. To address these issues we propose four extensions: (A) sample-weighted FedProx aggregation, giving a 2.6\% error reduction under data-size heterogeneity; (B) model-weight similarity clustering, which removes raw telemetry sharing entirely; (C) an SLA-aware asymmetric quantile loss ($\tau=0.8$), reducing SLA violations by 31\%; and (D) a clipping-and-noise differentially private variant (DP-FedProx) with an empirically characterised privacy--accuracy trade-off. All reported figures are traceable to the accompanying repository.

\medskip
\noindent\textbf{Keywords:} network slicing, RAN, federated learning, FedProx, traffic forecasting, differential privacy, reproducibility.
\end{abstract}

% ==================================================================
\section{Introduction and Motivation}
% ==================================================================
The advent of 5G and emerging 6G radio access networks (RANs) is underpinned by \textit{network slicing}, in which multiple virtualized, mutually isolated logical networks are dynamically multiplexed over a shared physical infrastructure. Physical base stations (BS) owned and operated by an Infrastructure Provider (InP) are leased by diverse Mobile Virtual Network Operators (MVNOs) and third-party vertical tenants to host application slices with widely different Quality of Service (QoS) and Service Level Agreement (SLA) profiles, such as enhanced Mobile Broadband (eMBB), Ultra-Reliable Low-Latency Communication (URLLC), and massive Machine-Type Communication (mMTC).

To maximize spectral efficiency and prevent costly SLA degradation, the InP must perform proactive radio resource pre-allocation (for example, Physical Resource Block slicing and virtual network function scheduling) on time scales of minutes. This relies on accurate spatio-temporal traffic forecasting~\cite{bega2019deepcog}. Classical deep learning approaches, however, face three practical impediments:
\begin{enumerate}
    \item \textbf{Privacy and commercial sovereignty.} Centralized aggregation requires MVNOs to transmit raw, user-identifiable packet counters and telemetry to a central InP repository. This conflicts with privacy legislation (e.g., GDPR) and exposes proprietary business information (subscriber densities, surge patterns, churn) to competitors.
    \item \textbf{Data scarcity and overfitting.} Isolated local training at individual base stations (Decentralized LSTM, DLSTM) overfits and generalizes poorly during non-stationary traffic bursts because of limited local history.
    \item \textbf{Backhaul overhead.} Continuous streaming of high-resolution multi-slice telemetry from hundreds of urban base stations saturates constrained backhaul links.
\end{enumerate}

To reconcile privacy preservation with forecasting accuracy, Phyu et al.~\cite{phyu2023multi} proposed a federated learning (FL) framework based on FedAvg~\cite{mcmahan2017communication}, FedProx~\cite{li2020federated}, and intelligent clustering (IC-FPLSTM). Although promising, the original paper contains architectural gaps, algorithmic inconsistencies, and insufficiently supported claims. This report makes three contributions:
\begin{itemize}
    \item \textbf{Faithful multi-seed replication.} All baselines and federated schemes are implemented and evaluated across fixed random seeds ($42$, $1337$, $2024$).
    \item \textbf{Ten-point critical audit.} Mathematical arguments, verbatim text comparisons, and empirical demonstrations expose flaws in Algorithms~1, 2, and 4, in the evaluation protocol, and in the privacy assumptions.
    \item \textbf{Four proposed improvements.} Concrete extensions that fix the theoretical mismatches, remove telemetry leakage, optimize an asymmetric SLA-oriented loss, and add differential-privacy mechanisms.
\end{itemize}

% ==================================================================
\section{System Model and Problem Formulation}
% ==================================================================

\subsection{RAN Slicing and Multi-Tenant Topology}
We model a metropolitan RAN deployment comprising a set $\mathcal{B} = \{1, 2, \dots, N\}$ of $N$ physical base stations. Each base station serves a common set of application slices $\mathcal{S} = \{1, 2, \dots, S\}$ (e.g., Facebook, YouTube, Google, Instagram). The traffic demand of slice $s \in \mathcal{S}$ at base station $b \in \mathcal{B}$ is discretized into uniform temporal bins of duration $\Delta t = 10$ minutes. Let $x_b^s(t) \in \mathbb{R}$ denote the aggregated traffic volume (e.g., in bytes or packets) served by base station $b$ for slice $s$ during the $t$-th interval.

Each base station maintains a local observation history of $H$ contiguous past time steps:
\begin{equation}
    \mathbf{X}_b^s(t) = \left[ x_b^s(t - H + 1), x_b^s(t - H + 2), \dots, x_b^s(t) \right]^{\top} \in \mathbb{R}^{H \times 1}.
\end{equation}
The forecasting objective is to predict the traffic volume over a horizon $W$ ($W=1$ in the paper's primary setup):
\begin{equation}
    \mathbf{Y}_b^s(t) = \left[ x_b^s(t + 1), \dots, x_b^s(t + W) \right]^{\top} \in \mathbb{R}^{W \times 1}.
\end{equation}

\subsection{Global Federated Learning Objective}
Each base station $b$ acts as a decentralized client holding a local dataset $\mathcal{D}_b = \{(\mathbf{X}_{b,i}, \mathbf{Y}_{b,i})\}_{i=1}^{D_b}$ with $D_b = |\mathcal{D}_b|$ samples. The total number of training instances is $D = \sum_{b=1}^N D_b$. The global objective, given as Eq.~(1) of the original paper, is a sample-weighted empirical risk minimization:
\begin{equation}
    \min_{\mathbf{w}} F(\mathbf{w}) \triangleq \sum_{b=1}^N \frac{D_b}{D} F_b(\mathbf{w}),
    \label{eq:global_obj}
\end{equation}
where $F_b(\mathbf{w}) = \frac{1}{D_b} \sum_{i=1}^{D_b} \ell\big(f(\mathbf{X}_{b,i}; \mathbf{w}), \mathbf{Y}_{b,i}\big)$ is the local empirical loss at base station $b$ and $\ell(\cdot)$ is the Mean Squared Error (MSE) criterion.

\subsection{LSTM Parameter Count and Complexity}
Each client deploys a single-layer LSTM followed by a fully connected linear projection. An LSTM cell with input dimension $d_{\text{in}}$ and hidden dimension $h$ contains four gating mechanisms (forget gate $f_t$, input gate $i_t$, candidate state $\tilde{c}_t$, output gate $o_t$). Each gate has an input weight matrix $\mathbf{W}_{*x} \in \mathbb{R}^{h \times d_{\text{in}}}$, a recurrent weight matrix $\mathbf{W}_{*h} \in \mathbb{R}^{h \times h}$, and a bias vector $\mathbf{b}_* \in \mathbb{R}^h$, so that
\begin{equation}
    P_{\text{cell}} = 4 \left( h \, d_{\text{in}} + h^2 + h \right).
\end{equation}
The linear layer maps the final hidden state $\mathbf{h}_H \in \mathbb{R}^h$ to the prediction window:
\begin{equation}
    P_{\text{linear}} = (h + 1)\, W.
\end{equation}
Taking $d_{\text{in}} = H$, the total trainable parameter count given as Eq.~(2) of the paper is
\begin{equation}
    P = 4 \left( h^2 + h H + h \right) + (h + 1)\, W.
    \label{eq:paper_lstm_params}
\end{equation}
For $h = 64$, $H = 5$, and $W = 1$, Eq.~\eqref{eq:paper_lstm_params} gives $P = 17{,}985$ parameters, i.e.\ a serialized payload of $\psi = 4P = 71{,}940$ bytes ($\approx 70.25$ KiB) in 32-bit floating-point precision.

\subsection{Communication Cost Model}
The cumulative communication overhead over $C$ federated rounds, with client participation fraction $\Omega \in (0, 1]$, is
\begin{equation}
    \text{Comm}_{\text{total}} = 2\, C\, N\, \Omega\, \psi \quad [\text{bytes}],
    \label{eq:comm_cost}
\end{equation}
where the factor of 2 accounts for the bidirectional transfer (downlink global broadcast and uplink local update).

\subsection{SLA Violation Rate}
To quantify the operational penalty of under-prediction, we define the SLA violation rate for an over-provisioning margin $\gamma \ge 0$:
\begin{equation}
    \text{SLA}_{\text{violation}}(\gamma) \triangleq \frac{1}{|\mathcal{D}_{\text{test}}|} \sum_{i=1}^{|\mathcal{D}_{\text{test}}|} \mathbb{I}\!\left( \hat{y}_i (1 + \gamma) < y_i \right),
    \label{eq:sla_violation}
\end{equation}
where $\mathbb{I}(\cdot)$ is the indicator function. A violation occurs whenever the allocated capacity $\hat{y}_i(1+\gamma)$ fails to meet the true demand $y_i$, causing packet drops and QoS penalties. The corresponding compliance rate is $1 - \text{SLA}_{\text{violation}}(\gamma)$.

% ==================================================================
\section{Methodology and Algorithms}
% ==================================================================

\subsection{Algorithm 1: Federated Learning Training Loop}
The coordination procedure follows an iterative master--worker protocol. At round $c=0$ the server initializes the global weights $\mathbf{w}^0$. In each round $c \in \{0, 1, \dots, C-1\}$:
\begin{enumerate}
    \item The server samples a random subset of base stations $\mathcal{B}_c \subseteq \mathcal{B}$ with $|\mathcal{B}_c| = \max(1, \lfloor \Omega N \rfloor)$.
    \item The server broadcasts $\mathbf{w}^c$ to all sampled clients.
    \item Each client $b \in \mathcal{B}_c$ runs local training (Algorithm~2) on $\mathcal{D}_b$ starting from $\mathbf{w}^c$ and returns $\mathbf{w}_b^{c+1}$.
    \item The server aggregates the received models into the new global model $\mathbf{w}^{c+1}$.
\end{enumerate}

\subsection{Algorithm 2: FedProx Local Descent}
To limit client drift under non-IID traffic distributions, FedProx~\cite{li2020federated} adds a proximal term to the local objective:
\begin{equation}
    h_b(\mathbf{w}; \mathbf{w}^c) \triangleq F_b(\mathbf{w}) + \frac{\mu}{2} \left\| \mathbf{w} - \mathbf{w}^c \right\|_2^2,
    \label{eq:fedprox_obj}
\end{equation}
where $\mu \ge 0$ controls the proximal pull toward the global anchor $\mathbf{w}^c$. In true FedProx, each local stochastic gradient step is
\begin{equation}
    \mathbf{w} \leftarrow \mathbf{w} - \eta \Big( \nabla F_b(\mathbf{w}; \xi) + \mu (\mathbf{w} - \mathbf{w}^c) \Big),
    \label{eq:true_fedprox_step}
\end{equation}
with learning rate $\eta$ and mini-batch $\xi$. Section~\ref{sec:critique} audits the paper's literal Algorithm~2, which decouples SGD from the proximal minimization.

\subsection{Algorithm 3: Intelligent Clustering (IC-FPLSTM)}
Instead of one global model across heterogeneous base stations, Algorithm~3 partitions them into $K$ clusters:
\begin{enumerate}
    \item \textbf{Feature extraction.} Base stations are described either by
    \begin{itemize}
        \item \textit{Traffic trend:} a 144-dimensional vector of the average 24-hour diurnal profile (144 ten-minute intervals), or
        \item \textit{Geographic coordinates:} a 2-dimensional vector $(\text{latitude}, \text{longitude})$.
    \end{itemize}
    \item \textbf{Dimensionality reduction.} Principal Component Analysis (PCA) projects the trend vectors onto 2 principal components, retaining the dominant variance.
    \item \textbf{Cluster count selection.} The Elbow method evaluates the sum of squared distances for $K \in \{2, \dots, 10\}$ and picks the inflection point.
    \item \textbf{Clustering.} $K$-means partitions the reduced representation into disjoint clusters $\mathcal{C}_1, \dots, \mathcal{C}_K$.
    \item \textbf{Independent federations.} A separate federated model is trained per cluster, aggregating only among its members.
\end{enumerate}

\subsection{Algorithm 4: Random Clustering (RC-FPLSTM)}
As an empirical baseline, Algorithm~4 partitions the $N$ base stations into $K$ random subsets of target size $\lfloor N / K \rfloor$. As detailed in Section~\ref{sec:critique}, the original Algorithm~4 contains a dynamic set-mutation bug, which we identify, analyse, and correct.

% ==================================================================
\section{Experimental Setup}
% ==================================================================

\subsection{Dataset Calibration}
The proprietary Orange Poitiers dataset used in the original paper is not publicly available. We therefore built a cellular traffic generator calibrated to the statistical parameters documented in Section~IV of the paper:
\begin{itemize}
    \item \textbf{Topology:} $N = 57$ base stations on a $10 \times 10\ \text{km}^2$ metropolitan grid.
    \item \textbf{Slices:} $|\mathcal{S}| = 4$ application slices:
    \begin{itemize}
        \item \textit{Facebook:} moderate volume, daytime burstiness, log-normal noise;
        \item \textit{YouTube:} heavy bandwidth consumption, pronounced evening peak (19:00--23:00);
        \item \textit{Google:} working-hours activity (09:00--17:00), low variance;
        \item \textit{Instagram:} erratic peaks, multi-modal diurnal spikes.
    \end{itemize}
    \item \textbf{Granularity and duration:} 10-minute intervals over 10 contiguous days (1,440 time steps per slice per base station; 228 time series in total).
    \item \textbf{Partitioning:} chronological split into 7 days training (70\%), 1 day validation (10\%), and 2 days out-of-sample testing (20\%), with no future look-ahead.
\end{itemize}

\subsection{Hyperparameter Configuration}
Following the paper: observation window $H = 5$, prediction window $W = 1$, hidden units $h = 64$, batch size 64, learning rate $\eta = 0.001$ (Adam/SGD), local epochs $E = 1$, communication rounds $C = 6$ to $20$, and participation fraction $\Omega \in [0.1, 1.0]$. All experiments use three fixed random seeds ($42$, $1337$, $2024$), and results are reported as mean $\pm$ standard deviation over seeds.

% ==================================================================
\section{Experimental Reproduction and Results}
% ==================================================================

\subsection{Baselines: Centralized LSTM vs.\ DLSTM}
Table~\ref{tab:fl_reproduction} summarizes the multi-seed results for all baselines and federated schemes.

\begin{table}[htbp]
\centering
\caption{Multi-seed test RMSE (mean $\pm$ std over seeds $42$, $1337$, $2024$) for baselines and federated schemes. The last column gives the change relative to DLSTM, computed from the reported means.}
\label{tab:fl_reproduction}
\small
\begin{tabular}{@{}llcc@{}}
\toprule
\textbf{Method} & \textbf{Setting} & \textbf{Test RMSE} & \textbf{vs.\ DLSTM} \\
\midrule
Centralized LSTM          & Raw data pooled at InP      & $0.0712 \pm 0.0039$ & $-65.1\%$ \\
DLSTM                     & Isolated local training     & $0.2039 \pm 0.0075$ & --- \\
FALSTM                    & FedAvg ($\mu = 0$)          & $0.1748 \pm 0.0022$ & $-14.3\%$ \\
FPLSTM                    & True FedProx ($\mu = 1.0$)  & $0.1757 \pm 0.0014$ & $-13.8\%$ \\
FPLSTM (literal Alg.~2)   & Paper's Algorithm 2 as written & $0.1850$         & $-9.3\%$ \\
\bottomrule
\end{tabular}
\end{table}

The centralized LSTM baseline sets an empirical lower bound on error, with an average test RMSE of $0.0712 \pm 0.0039$ across all slices. Isolated DLSTM models suffer from data starvation and reach an average test RMSE of $0.2039 \pm 0.0075$, which is $186\%$ higher than the centralized model. DLSTM performs particularly poorly on bursty slices (Facebook: $0.2602$; YouTube: $0.2334$).

\subsection{FL Reproduction: FALSTM vs.\ FPLSTM}
Federated Averaging (FALSTM, $\mu=0$) reaches a test RMSE of $0.1748 \pm 0.0022$, outperforming isolated DLSTM by $14.3\%$ while keeping raw telemetry on the local base stations. True FedProx (FPLSTM, $\mu=1.0$) reaches $0.1757 \pm 0.0014$ and has the lowest standard deviation across seeds, which is consistent with proximal regularization stabilizing inter-client divergence under non-IID conditions.

\subsection{Proximal Regularization Sweep ($\mu$)}
Sweeping $\mu \in \{0.001, 0.01, 0.1, 0.5, 1.0\}$ at $\Omega = 0.5$ confirms the paper's finding that larger proximal penalties reduce generalization error: RMSE falls from $0.1804$ at $\mu=0.001$ to $0.1743$ at $\mu=1.0$. A strong proximal pull anchors client updates and prevents over-specialization on noisy local samples.

\subsection{Client Participation Rate ($\Omega$)}
For $\Omega \in \{0.1, 0.3, 0.5, 0.7, 1.0\}$, the RMSE drops sharply from $0.2517$ at $\Omega=0.1$ to $0.1810$ at $\Omega=0.3$ and then plateaus ($0.1743$ at $\Omega=1.0$). Cumulative communication overhead grows linearly from $17.97$~MB to $185.04$~MB. Hence $\Omega = 0.5$ is a sensible operating point, reaching about 98\% of full-participation accuracy at half the radio bandwidth overhead.

\subsection{Clustered Federated Learning}
Table~\ref{tab:clustering_results} compares clustering strategies with $K=3$.

\begin{table}[htbp]
\centering
\caption{Multi-seed test RMSE of clustered FL strategies ($K=3$). The last column gives the relative improvement over corrected random clustering (RC-Fixed).}
\label{tab:clustering_results}
\small
\begin{tabular}{@{}llcc@{}}
\toprule
\textbf{Method} & \textbf{Clustering feature} & \textbf{Test RMSE} & \textbf{Gain vs.\ RC-Fixed} \\
\midrule
IC-FPLSTM  & 24-h traffic trend (PCA) & $0.2149 \pm 0.0052$ & $1.24\%$ \\
IC-FPLSTM  & Latitude / longitude     & $0.2116 \pm 0.0120$ & $2.76\%$ \\
RC-FPLSTM  & Random, corrected (Fixed)& $0.2176 \pm 0.0033$ & --- \\
\bottomrule
\end{tabular}
\end{table}

IC-FPLSTM (Trend) achieves an average RMSE of $0.2149 \pm 0.0052$ and IC-FPLSTM (Lat/Long) achieves $0.2116 \pm 0.0120$, while corrected random clustering (RC-FPLSTM Fixed) achieves $0.2176 \pm 0.0033$. The differences are small relative to the seed-to-seed standard deviations, especially for Lat/Long, whose gain is below its own standard deviation.

% ==================================================================
\section{Critical Audit: Ten Architectural Flaws}
\label{sec:critique}
% ==================================================================
Our systematic audit identified ten significant weaknesses, summarized in Table~\ref{tab:audit_matrix} and discussed in the subsections that follow.

{\footnotesize
\begin{longtable}{@{}p{0.95cm}p{1.9cm}p{6.0cm}p{5.2cm}@{}}
\caption{Ten-point critical audit of Phyu et al.~\cite{phyu2023multi}.}
\label{tab:audit_matrix}\\
\toprule
\textbf{ID} & \textbf{Category} & \textbf{Discovered issue and evidence} & \textbf{Empirical verdict / impact} \\
\midrule
\endfirsthead
\toprule
\textbf{ID} & \textbf{Category} & \textbf{Discovered issue and evidence} & \textbf{Empirical verdict / impact} \\
\midrule
\endhead
\midrule
\multicolumn{4}{r}{\emph{Continued on next page}}\\
\endfoot
\bottomrule
\endlastfoot
\textbf{A-01} & Algorithm & Alg.~2 runs vanilla SGD without the proximal term and applies a post-hoc $\arg\min$ in Line~9. & Literal Alg.~2 yields RMSE $0.1850$ vs.\ $0.1757$ for true FedProx ($+5.3\%$ error). \\[3pt]
\textbf{A-02} & Mathematics & The paper's Eq.~(1) defines the weighted sum $\sum_b \tfrac{D_b}{D} F_b$, but Alg.~1 Line~7 computes the unweighted mean $\tfrac{1}{N}\sum_b \mathbf{w}_b$. & Biased convergence under data heterogeneity; sample weighting reduces error by $2.6\%$. \\[3pt]
\textbf{A-03} & Algorithm & Alg.~4 mutates the set, $\mathcal{B}_s \leftarrow \mathcal{B}_s \setminus \{b\}$, inside the while loop, shrinking $|\mathcal{B}_s|/K$ dynamically. & Early loop termination; up to $33\%$ of stations end up in the surplus cluster. \\[3pt]
\textbf{A-04} & Notation & Sec.~VI-E defines $H$ as prediction and $W$ as observation window, swapping the definitions of Sec.~IV-B. & Notation inconsistency; hinders direct replication without reconciliation. \\[3pt]
\textbf{A-05} & Evaluation & Single city (Poitiers), 10 days, no repeated seeds, no confidence intervals, hyperparameters tuned on the test set. & Weak empirical robustness; results may reflect initialization artifacts. \\[3pt]
\textbf{A-06} & Baseline & DeepCog~\cite{bega2019deepcog} is trained on an asymmetric capacity loss but evaluated only with symmetric RMSE; communication model is flawed. & Unfair baseline comparison; ignores the long-term communication crossover. \\[3pt]
\textbf{A-07} & Overclaim & The abstract claims IC-FPLSTM is clearly superior, whereas Table~V shows a gain of only $0.4\%$--$1.5\%$. & Our multi-seed run shows a gain of only $1.24\%$ over balanced random clustering. \\[3pt]
\textbf{A-08} & Privacy & IC-Trend requires clients to send raw 24-hour traffic vectors; no formal DP guarantee. & Violates the FL privacy premise; exposes fine-grained tenant profiles. \\[3pt]
\textbf{A-09} & Realism & A single global model is forced across heterogeneous slices; $K=3$ is fixed arbitrarily; clustering CPU time is omitted. & Impractical for independent MVNOs; hides the overhead of PCA and Elbow. \\[3pt]
\textbf{A-10} & Benchmarking & DLSTM vs.\ FPLSTM CPU time (54 min vs.\ 8.5 min) reflects 11,400 vs.\ 20 training passes. & Artifact of epoch selection, not inherent algorithmic efficiency. \\
\end{longtable}
}

\subsection{Point 1: Algorithm 2 vs.\ True FedProx}
The paper's Algorithm~2 states:
\begin{quote}
\textit{``Line 8: Update weights using Adam optimizer; Line 9: $\mathbf{w}_b^{c+1} = \arg\min_{\mathbf{w}} \big(F_b(\mathbf{w}) + \tfrac{\mu}{2}\|\mathbf{w} - \mathbf{w}^c\|^2\big)$.''}
\end{quote}
This is internally inconsistent. Line~8 performs standard local SGD for $E$ epochs without any proximal regularization, and Line~9 then invokes a post-hoc $\arg\min$ that has no counterpart in standard backpropagation. True FedProx adds the proximal gradient $\mu(\mathbf{w} - \mathbf{w}^c)$ to every mini-batch step, as in Eq.~\eqref{eq:true_fedprox_step}. Empirically, the literal Algorithm~2 yields an RMSE of $0.1850$, which is $5.3\%$ worse than true FedProx ($0.1757$).

\subsection{Point 2: Unweighted Aggregation Mismatch}
The paper's Eq.~(1) defines the global risk as the sample-weighted sum in Eq.~\eqref{eq:global_obj}, yet Algorithm~1 Line~7 performs unweighted averaging:
\begin{equation}
    \mathbf{w}^{c+1} = \frac{1}{|\mathcal{B}_c|} \sum_{b \in \mathcal{B}_c} \mathbf{w}_b^{c+1}.
\end{equation}
When base stations hold unequal amounts of data $D_b$ (for example, due to staggered deployment or missing telemetry), unweighted aggregation over-weights data-scarce stations and biases the global update.

\subsection{Point 3: Algorithm 4 Set-Mutation Bug}
Lines 3--8 of Algorithm~4 specify the loop shown in Algorithm~\ref{alg:alg4_original}.

\begin{algorithm}[htbp]
\caption{Cluster filling loop of the original Algorithm 4 (as specified in the paper)}
\label{alg:alg4_original}
\begin{algorithmic}[1]
\While{$|c_k| < |\mathcal{B}_s| / K$}
    \State Sample $b \in \mathcal{B}_s$
    \State $c_k \leftarrow c_k \cup \{b\}$
    \State $\mathcal{B}_s \leftarrow \mathcal{B}_s \setminus \{b\}$
\EndWhile
\end{algorithmic}
\end{algorithm}

Because $\mathcal{B}_s$ is pruned inside the loop, the threshold $|\mathcal{B}_s|/K$ decreases after every assignment. Cluster $c_1$ therefore stops after acquiring only about $N/(K+1)$ elements, since the loop condition becomes $|c_1| < (N - |c_1|)/K$. The remaining stations fall into the surplus branch, producing strongly imbalanced clusters (e.g., sizes $12, 16, 29$ instead of $19, 19, 19$).

\subsection{Point 7: Overstated Clustering Superiority}
The abstract of the original paper claims that IC-FPLSTM clearly outperforms random clustering. However, the paper's own Table~V shows that IC-FPLSTM improves RMSE over RC-FPLSTM by only $0.4\%$ to $1.5\%$. In our multi-seed reproduction (Table~\ref{tab:clustering_results}), IC-Trend achieves $0.2149$ vs.\ $0.2176$ for RC-Fixed, a gain of only $1.24\%$.

\subsection{Point 8: Telemetry Leakage in Intelligent Clustering}
The core premise of federated learning is that raw client telemetry never leaves the base station. Yet Algorithm~3 requires every base station to send its full 24-hour traffic trend vector to the central InP coordinator for PCA and $K$-means. This exposes fine-grained tenant traffic profiles and undermines the privacy rationale of FL.

% ==================================================================
\section{Proposed Improvements and Evaluations}
% ==================================================================

\subsection{Improvement A: Sample-Weighted FedProx}
To resolve Audit Point~2, we use exact sample-weighted aggregation:
\begin{equation}
    \mathbf{w}^{c+1} = \sum_{b \in \mathcal{B}_c} \frac{D_b}{\sum_{j \in \mathcal{B}_c} D_j}\, \mathbf{w}_b^{c+1}.
\end{equation}
To benchmark it, we simulated dataset-size heterogeneity by varying each base station's observation history between 50\% and 100\% of the full window. Sample-weighted aggregation reaches a test RMSE of $0.2242 \pm 0.0018$, compared with $0.2303 \pm 0.0021$ for unweighted averaging. This is a $2.6\%$ error reduction, with gains above $11.7\%$ on individual seeds.

\subsection{Improvement B: Model-Weight Similarity Clustering}
To remove the leak identified in Audit Point~8, we propose \textit{model-weight similarity clustering}. Clients run a single local warm-up round from a common initial model $\mathbf{w}^0$. Instead of raw telemetry, each client uploads its model update $\Delta\mathbf{w}_b = \mathbf{w}_b^1 - \mathbf{w}^0$. The server computes pairwise cosine similarity between update vectors:
\begin{equation}
    \text{Sim}(b, j) = \frac{\Delta\mathbf{w}_b \cdot \Delta\mathbf{w}_j}{\left\|\Delta\mathbf{w}_b\right\|_2 \left\|\Delta\mathbf{w}_j\right\|_2},
\end{equation}
and clusters clients with $K$-means in parameter space. This reaches an RMSE of $0.2855$ without any raw traffic telemetry being sent to the server. The error is higher than that of trend-based clustering ($0.2149$), so the privacy gain comes with an accuracy cost that deployments must weigh.

\subsection{Improvement C: SLA-Aware Asymmetric Quantile Loss}
MSE penalizes over- and under-prediction symmetrically. In production slicing, under-prediction causes immediate SLA breaches and penalties, whereas over-prediction only wastes some radio resources. We therefore use the asymmetric pinball loss with target quantile $\tau = 0.8$:
\begin{equation}
    \mathcal{L}_{\tau}(y, \hat{y}) = \max\!\Big( \tau (y - \hat{y}),\; (1 - \tau)(\hat{y} - y) \Big).
\end{equation}
This loss penalizes under-allocation $\tau/(1-\tau) = 4\times$ more than over-allocation. At $\gamma=0$, SLA compliance rises from $78.23\%$ under MSE to $85.04\%$ under pinball loss, so the SLA violation rate on the test partition falls from $21.77\%$ to $14.96\%$, a \textbf{31\% relative reduction}. The RMSE changes by only $0.0026$ ($0.3356$ vs.\ $0.3382$).

\subsection{Improvement D: Differentially Private FedProx (DP-FedProx)}
To mitigate membership-inference and reconstruction attacks (Audit Point~8), we add differential-privacy mechanisms~\cite{dwork2014algorithmic} to the federated loop. In each round, local updates $\Delta\mathbf{w}_b = \mathbf{w}_b^{c+1} - \mathbf{w}^c$ are clipped to a maximum $\ell_2$ norm $C_{\text{clip}} = 1.0$:
\begin{equation}
    \Delta\tilde{\mathbf{w}}_b = \Delta\mathbf{w}_b \cdot \min\!\left(1, \frac{C_{\text{clip}}}{\left\|\Delta\mathbf{w}_b\right\|_2}\right).
\end{equation}
Calibrated Gaussian noise $\mathcal{N}(0, \sigma^2 \mathbf{I})$ is then added before aggregation. We evaluated the privacy--accuracy trade-off for $\sigma \in \{0.0, 0.005, 0.01, 0.05, 0.1\}$. At $\sigma = 0.005$ the model achieves an RMSE of $0.1972$, retaining high predictive fidelity under noise injection. For $\sigma \ge 0.05$ the RMSE rises to $0.2015$, which marks the practical limit of the privacy--accuracy trade-off in our setting. A formal $(\varepsilon, \delta)$ accounting over the training rounds is left as future work.

\begin{table}[htbp]
\centering
\caption{Summary of the four proposed improvements and their measured effect.}
\label{tab:improvements_summary}
\small
\begin{tabularx}{\textwidth}{@{}l l Y Y@{}}
\toprule
\textbf{ID} & \textbf{Improvement} & \textbf{Addresses} & \textbf{Measured result} \\
\midrule
A & Sample-weighted FedProx & Audit A-02 (unweighted aggregation) & RMSE $0.2303 \pm 0.0021 \rightarrow 0.2242 \pm 0.0018$ ($-2.6\%$) under data-size heterogeneity \\[3pt]
B & Model-weight similarity clustering & Audit A-08 (telemetry leakage) & No raw telemetry shared; RMSE $0.2855$ (vs.\ $0.2149$ for IC-Trend) \\[3pt]
C & SLA-aware pinball loss ($\tau = 0.8$) & SLA-agnostic MSE objective & SLA violations $21.77\% \rightarrow 14.96\%$ ($-31\%$ relative) at $\gamma=0$ \\[3pt]
D & DP-FedProx (clipping $C_{\text{clip}}=1.0$ + Gaussian noise) & Audit A-08 (inference attacks) & RMSE $0.1972$ at $\sigma=0.005$; $0.2015$ for $\sigma \ge 0.05$ \\
\bottomrule
\end{tabularx}
\end{table}

% ==================================================================
\section{Individual Contributions}
% ==================================================================
In accordance with the course rubric, responsibilities were distributed as follows:

\begin{table}[htbp]
\centering
\small
\caption{Team roles and individual contributions.}
\label{tab:contributions}
\begin{tabularx}{\textwidth}{@{}l l Y@{}}
\toprule
\textbf{Member} & \textbf{Role} & \textbf{Contribution} \\
\midrule
Ankush R (CS24B1062) & Paper Analyst & Formalized the 5G RAN slicing problem; audited Table~II and Eqs.~(1)--(2); cataloged 12 paper assumptions; built the claims-vs-evidence matrix. \\[3pt]
Vihaan Gupta (CS24B1063) & Reproduction Lead & Designed the Poitiers-calibrated traffic generator; implemented the preprocessing and sliding-window pipeline; built the DLSTM and Centralized baselines and the core FALSTM and FPLSTM simulators. \\[3pt]
Piyush Mishra (CS24B1048) & Scalability Lead & Implemented IC-FPLSTM (PCA + Elbow + $K$-means) and RC-FPLSTM; found and fixed the Algorithm~4 set-mutation bug; profiled communication costs; ran the multi-seed scalability benchmarks. \\[3pt]
Narayana Murthy S (CS24B1047) & Extension Lead & Formulated the ten-point audit matrix; implemented the four extensions (Sample-Weighted FedProx, Model-Weight Clustering, SLA Pinball Loss, DP-FedProx); conducted the ablation studies. \\
\bottomrule
\end{tabularx}
\end{table}

% ==================================================================
\section{Conclusion}
% ==================================================================
This report presented an independent, end-to-end replication and critical audit of multi-slice federated traffic forecasting. We confirm that FedProx and FedAvg clearly outperform isolated local training (a $14.3\%$ RMSE reduction), but the audit also found errors in Algorithms~2 and~4, a mismatch between the stated objective and the aggregation rule, only marginal benefits from intelligent clustering, and significant privacy weaknesses. The four proposed extensions (sample-weighted aggregation, model-weight clustering, SLA-aware pinball loss, and differentially private updates) address these issues and yield measurable gains in accuracy, SLA compliance, and privacy. Remaining limitations include the synthetic, calibrated dataset, the accuracy cost of model-weight clustering, and the absence of a formal $(\varepsilon,\delta)$ privacy accounting for DP-FedProx. These are natural directions for future work.

% ==================================================================
% References
% ==================================================================
\begin{thebibliography}{10}

\bibitem{phyu2023multi}
H.~P. Phyu, R.~Stanica, and D.~Naboulsi, ``Multi-slice privacy-aware traffic forecasting at {RAN} level: A scalable federated-learning approach,'' \emph{IEEE Trans. Netw. Service Manag.}, vol.~20, no.~4, pp.~4621--4634, Dec. 2023.

\bibitem{li2020federated}
T.~Li, A.~K. Sahu, M.~Zaheer, M.~Sanjabi, A.~Talwalkar, and V.~Smith, ``Federated optimization in heterogeneous networks,'' in \emph{Proc. MLSys}, 2020, pp.~429--450.

\bibitem{mcmahan2017communication}
B.~McMahan, E.~Moore, D.~Ramage, S.~Hampson, and B.~A. y~Arcas, ``Communication-efficient learning of deep networks from decentralized data,'' in \emph{Proc. AISTATS}, 2017, pp.~1273--1282.

\bibitem{bega2019deepcog}
D.~Bega, M.~Gramaglia, M.~Fiore, A.~Banchs, and X.~Costa-P{\'e}rez, ``{DeepCog}: Cognitive network management in sliced 5{G} networks with deep learning,'' in \emph{Proc. IEEE INFOCOM}, 2019.

\bibitem{dwork2014algorithmic}
C.~Dwork and A.~Roth, ``The algorithmic foundations of differential privacy,'' \emph{Found. Trends Theor. Comput. Sci.}, vol.~9, no.~3--4, pp.~211--407, 2014.

\end{thebibliography}

\end{document}